\documentclass[11pt]{article}
\usepackage[letterpaper,margin=0.85in]{geometry}
\usepackage[T1]{fontenc}
\usepackage{lmodern,amsmath,booktabs,array,xcolor,fancyhdr}
\usepackage[colorlinks=true,linkcolor=blue,urlcolor=blue]{hyperref}
\definecolor{ink}{HTML}{17344A}
\pagestyle{fancy}\fancyhf{}
\fancyhead[L]{\small\textcolor{ink}{UpstreamDrift / Necromatcher}}
\fancyhead[R]{\small Local Research Golf Methods}
\fancyfoot[L]{\small Model-Contact Research; Qualification Unverified}
\fancyfoot[R]{\thepage}
\setlength{\headheight}{14pt}
\setlength{\parindent}{0pt}\setlength{\parskip}{7pt}
\newcommand{\code}[1]{\nolinkurl{#1}}
\title{\textcolor{ink}{Necromatcher Local Research Golf\\Authenticated Impact Admission and Explicit Simulation}}
\author{UpstreamDrift Technical Reference\\Epic \#11232; Integration Issue \#11493}
\date{4 October 2026 / Local Research Golf Supplement 1}
\begin{document}
\maketitle
\begin{center}\rule{0.9\linewidth}{0.6pt}\end{center}
\section*{Purpose and Research Boundary}
This procedure connects a saved, source-bound authored replay impact to the
existing local golf simulator. It is reusable across player records, models and
driving profiles; it does not create a player-specific physics implementation.
Model contact remains assumed, authored seconds remain uncalibrated historical
time, and contact, numerical and scientific qualification remain unverified.
Local delivery acceptance is an execution state, not a historical validity badge.

Tiger matching retains original frames 0--190, the conservative reviewed domain
$[0,191)$. The uncertain hand-release transition and later one-hand follow-through
remain excluded. Hogan retains $[0,750)$. Camera, anatomical dimensions, human
range of motion, continuous grip/contact, shaft flex and imaging distortion
remain separate historical qualification requirements.

\section*{Immutable Evidence Admission}
The public research bridge receives an explicit registered replay ID, completed
impact run ID and shot/session metadata. The canonical impact session repeats
parent, manifest, execution and exact artifact checks before a shot is prepared.
The portable bundle retains \code{request.json}, \code{result.json},
\code{impact-receipt.json} and \code{trajectory.json}. A successful rejected
research run may be used; incomplete, cancelled, foreign or altered evidence
cannot establish admission.

The companion receipt binds complete extraction state to exact trajectory bytes.
Result and request parents, geometry, selection and execution identities must
agree, as must the receipt's extraction metadata. The bridge retains hashes and
actual saved assumptions through a detached whitelist. Uncontrolled JSON fields
cannot override run identities or promote qualification. Original input files
and the retained original trajectory are not rewritten by the conversion.

\newpage
\section*{Launch Vectors and Declared Aim}
Use the complete saved post-impact ball velocity and angular velocity. Neither
clubhead scalar speed nor differences between plotted positions replace these
vectors. The existing post-impact bridge constructs the canonical shot envelope;
the existing spatial transformation owner applies the proper declared aim:
\[
\mathbf v_T=R_{TF}\mathbf v_F,\qquad
\boldsymbol\omega_T=R_{TF}\boldsymbol\omega_F,\qquad
R_{TF}^{\mathsf T}R_{TF}=I,\quad\det R_{TF}=+1.
\]
The source uses \code{flight_xfwd_yleft_zup}; the target is the declared target
local frame. Proper rotations preserve vector norms and handle angular velocity
without a reflection. Local two-dimensional contact offsets retain their own
club coordinates. The saved source vectors and byte hashes remain available
independently of this declared target transformation.

Model-run, trace digest, impact identity and selected authored sample time derive
from the authenticated result, rather than caller-supplied replacement lineage.
The sample is a declared hypothesis; no historical collision is detected by
preparation, arming or explicit submission.

\section*{Canonical Local Session Policy}
\begin{center}\small
\begin{tabular}{>{\raggedright\arraybackslash}p{0.28\linewidth}>{\raggedright\arraybackslash}p{0.63\linewidth}}\toprule
Boundary & Required Behavior\\\midrule
Ordinary Model Preparation & Retains the qualified-contact requirement.\\
Research Preparation & Explicit opt-in; actual connected local adapter; all three qualification axes unverified; nonempty evidence and model lineage.\\
Arm and Submit & Owned preparation and single-use token; recheck research destination before submission.\\
Destination Change & Invalidates pending preparation; research cannot use fake, relay or external destinations.\\
Trajectory Recall & Confirmed owned local receipt, exact shot identity and detached retained samples.\\\bottomrule
\end{tabular}
\end{center}
The local adapter executes its existing flight simulator outside the event loop.
It computes a new flight using its own settings and environment, rather than
claiming byte identity with the retained original impact trajectory. Saved
original assumptions remain visible, but are not silently adopted as new local
environment settings. Native avatar, course feedback and aim-control support
remain unsupported by this destination.

\newpage
\section*{Explicit Operator Sequence}
Recall a registered authored replay and verified completed impact result, then
choose \textbf{Open Local Research Simulation}. Review saved IDs, hashes,
assumptions and authored sample time. Choose \textbf{Connect}, then
\textbf{Prepare Research Shot}, \textbf{Arm} and explicitly \textbf{Submit}.
Preparation does not arm or send automatically. Disarm or cancel through actual
owned controls; source or destination replacement must retire stale preparation
before new controls become usable.

The local delivery receipt and retained sample positions/velocities are distinct
from scientific acceptance. Web visualization reuses the existing flight scene;
native review exposes retained values. Malformed research links do not select a
manual demo. Failed or stale responses cannot restore old tokens or silently
connect another destination.

\section*{API and Persistence}
The research preparation endpoint accepts explicit replay/run/shot/session IDs,
timestamp, proper aim and context revision. Expensive evidence authentication
runs outside the API event loop; session replacement during that operation
refuses preparation. Existing lifecycle endpoints retain their canonical owners.
After confirmed local delivery, trajectory retrieval returns new samples with
the source research context. Reused shot IDs cannot replace prior context.

Registered replays and original impact results persist in the canonical Library.
Local golf sessions and their current trajectories are session data; restart
recall of a local simulation is not claimed. After restarting, recall the saved
source result and explicitly create another local research simulation.

\section*{Verification Protocol and Remaining Work}
Meaningful failing tests precede implementation. Verify full asymmetric launch
and spin vectors, nonidentity proper aim, qualification-promotion refusal,
foreign/stale/altered bundles, one-shot lifecycle, cancellation and responsive
background work. Ordinary service/GUI doubles establish boundary behavior only.
Native acceptance uses a fresh temporary canonical Library and actual MuJoCo
replay, clean-worker impact, Rust flight and real local submission. Exact
implementation bytes, terminal process logs, retained samples and inspected
screenshots belong to the accompanying frozen review.

This integration does not qualify a Tiger/Hogan match, establish physical source
time, animate a native avatar or prove external/course delivery. Historical
matching and final overall integration remain active. The separate engineering
manual inventory and approval gates remain in force.

\section*{References and Public Owners}
\href{https://github.com/D-sorganization/UpstreamDrift/issues/11493}{Issue \#11493}
and \href{https://github.com/D-sorganization/UpstreamDrift/pull/11359}{PR \#11359}.
Public owners: the research impact bridge, \code{GolfSessionService},
\code{LocalReferenceAdapter}, post-impact launch conversion and the existing
spatial transformation and visualization providers.
\end{document}
